The ordered case provides the starting point. Immerman and Vardi showed
that first-order logic with a least fixed-point operator expresses exactly
the polynomial-time queries on finite structures supplied with a linear
order~\cite{immerman1982,vardi1982}.
For unordered structures, Chandra and Harel asked whether polynomial-time
queries admit an effective enumeration~\cite[Section~5, p.~118]{chandraHarel1982}.
Gurevich formulated the question in terms of a logic capturing polynomial
time and conjectured that no such logic exists~\cite[Section~7]{gurevich1988}.
Excluding CPT with counting as a candidate leaves this broader question
open.

Blass, Gurevich and Shelah introduced choiceless polynomial time to study
computation on unordered structures using hereditarily finite sets rather
than arbitrary choices of an ordering~\cite{bgs1999}. Their original
1997 preprint explicitly conjectured that the extension by counting is
still a proper fragment of polynomial time~\cite[1997 version,
Introduction, p.~3]{bgs1999}. Their later study of unordered computation
examined CFI instances and finite-field linear algebra as possible sources
of separation, carefully distinguishing representations with common and
independent row and column index sets~\cite[Sections~4 and~6]{bgs2002}.
The query here is a particular linear-consistency problem over
$\mathbb F_3$, defined without a promise on the input structure.

Shelah's early work claimed noncapture for counting extensions of
choiceless computation~\cite[Introduction and Section~4]{shelah2000}.
Later accounts continued to state the capture problem as open, from
Blass, Gurevich and Shelah~\cite[Section~7]{bgs2002} through the 2025
work of Dawar, Gr{\"a}del, Kullmann and Pago~\cite[Section~2.2]{dawarGraedelKullmannPago2025}.
We give a separate proof for the explicit query above; Shelah's claim is
not a theorem input to this argument.

The construction of Cai, F{\"u}rer and Immerman established the central
example of a polynomial-time graph query beyond fixed-point logic with
counting~\cite{cfi1992}. Its local gadgets conceal a global parity
obstruction from bounded-variable counting formulas. This is an
antecedent of our charge-based witnesses, not an identification of those
witnesses with the original CFI graphs. The passage from counting logic
to choiceless computation is substantial: Dawar, Richerby and Rossman
showed that the preordered CFI query obtained from ordered base graphs is
definable without counting, while even counting does not suffice if set
rank is restricted to $o(\log n/\log\log n)$, where $n$ is the input
size~\cite[Theorems~17 and~40,
2007 author version]{drr2008}. In particular, their lower bound for every
fixed rank cannot by itself yield a separation for the full formalism.
Pakusa, Schalth{\"o}fer and Selman extended the positive result to
connected preordered base graphs whose colour classes have logarithmic
size in the number of base vertices~\cite[Theorem~1]{pakusaSchalthoeferSelman2016}.
The preorder orders classes of base vertices, rather than all vertices
of the resulting CFI structure.

Positive algebraic results also go beyond the original CFI family.
Abu Zaid, Gr{\"a}del, Grohe and Pakusa obtained choiceless canonization,
that is, construction of canonical ordered copies, when a supplied preorder has bounded classes whose induced substructures
have Abelian automorphism groups~\cite[Corollary~19]{abuZaidGraedelGrohePakusa2014}.
Their procedure uses specially structured cyclic systems of equations,
not arbitrary unordered linear systems~\cite[Definition~6 and
Theorem~14]{abuZaidGraedelGrohePakusa2014}. Lichter and Schweitzer obtained
analogous canonization for graphs when these induced groups are dihedral
or cyclic~\cite[Theorem~1]{lichterSchweitzer2021}. These results show why
the lower-bound argument must exploit the particular action of its
nonabelian group, rather than noncommutativity alone.

Functional lower bounds expose different limits. Rossman proved that
CPT with counting cannot construct the set of hyperplanes of an input
finite vector space over a fixed finite
field~\cite[Theorem~6.1, author version]{rossman2010}.
Pago proved that no such program constructs, on every hypercube with
$N$ vertices, a total preorder with classes of size
$O(\log N)$~\cite[Theorem~2 and Corollary~3, full version]{pago2021}.
The latter statement also excludes encodings from which the preorder can
be recovered in CPT. It obstructs the route of first constructing a fine
preorder and then applying the preordered CFI algorithm. These are
functional obstructions, not nondefinability of a Boolean query.

Pago's symmetric-XOR-circuit approach relates restricted classes of
choiceless CFI algorithms to circuit
families~\cite[Theorems~2, 31 and~33, full version]{pago2023}.
The translation permits polynomially many circuit images under
automorphisms of the base graph; the circuit lower bound assumes an
invariant circuit and different bounds on connections between gates.
These hypotheses differ, so that argument does not establish
nondefinability for the entire restricted algorithm class.

The support-and-transfer architecture has older roots. Blass, Gurevich
and Shelah developed supports and representations by forms evaluated at
supporting tuples, called molecules~\cite[Sections~8--9, published
version]{bgs1999}. Dawar, Richerby and Rossman developed rank-sensitive
support bounds and a counting transfer for hereditarily supported
objects~\cite[Sections~7--8, 2007 author version]{drr2008}. Using the
group-relative homogeneity from Section~\ref{sec:base-counting}, we adapt
that representation method in Section~\ref{sec:hf-transfer} and prove the
required exact counting statements.
The support estimate in Section~\ref{sec:support} instead uses a larger
nonabelian group to obtain small supports for its central subgroup,
independently of HF rank. Its standard class-two group and
alternating-form ingredients, as well as the projection and counting-game
methods used for the atom comparison, are credited at their respective
proofs. The grid-specific estimate and the quantitative transfer are
proved here.

Finally, the connection to arbitrary programs uses the interpretation
characterization of Gr{\"a}del, Pakusa, Schalth{\"o}fer and
Kaiser~\cite[Theorem~1]{gkps2015}. We use the binary, two-dimensional
formulation of Grohe, Schweitzer and Wiebking~\cite[Section~6, Theorem~18,
arXiv version~1]{gsw2020}, including its generated-equivalence and
existential-representative quotient convention. This characterization is
an external input. Section~\ref{sec:coding} proves the further facts
needed here: a polynomial bound on the complete hereditary encoding of a
run, and exact formulas for its states whose variable bound is independent
of the number of stages. Thus neither a bound on HF rank nor a restriction
to a particular algorithmic construction is imposed on the programs
excluded by the theorem.
